\section{Управление учетной записью}\label{ACME.Account}

\subsection{Создание учетной записи}\label{ACME.Account.Create}

Клиент создает новую учетную запись, отправляя
\pdfexample{POST-запрос}{acme/account1.txt} на узел \code{newAccount}. Тело
запроса представляет собой незавершенный объект учетной записи 
(см.~\ref{ACME.Res.Account}) с необязательными полями \sstr{contact}, 
\sstr{termsOfServiceAgreeed}, \sstr{externalAccountBinding} и одним 
дополнительным полем:

\begin{itemize}
\item
\sstr{onlyReturnExisting} (необязательное, логическое значение)~---
если поле присутствует и в нем указано значение \code{true}, то сервер не 
должен создавать новую учетную запись. Это позволяет клиенту искать адрес уже 
существующей учетной записи по ее ключу (см.~\ref{ACME.Account.Find}).
\end{itemize}

Поля \sstr{contact}, \sstr{termsOfServiceAgreeed} и 
\sstr{externalAccountBinding} заполняются по правилам, установленным 
в~\ref{ACME.Res.Account}.

\if0
\begin{codeblock}
POST /acme/new-account HTTP/1.1
Host: example.com
Content-Type: application/jose+json

{
  "protected": base64url({
    "alg": "BIGNS128",
    "jwk": {...},
    "nonce": "6S8IqOGY7eL2lsGoTZYifg",
    "url": "https://example.com/acme/new-account"
  }),
  "payload": base64url({
    "termsOfServiceAgreed": true,
    "contact": [
      "mailto:cert-admin@example.org",
      "mailto:admin@example.org"
    ]
  }),
  "signature": "RZPOnYoPs1PhjszF...-nh6X1qtOFPB519I"
}
\end{codeblock}
\fi

Сервер должен игнорировать любые значения, указанные в поле \sstr{orders} 
объекта учетной записи, отправленного клиентом, а также любые другие поля, 
которые сервер не распознает. 
%
Если в будущем появится дополнительное поле, то в его спецификации должно быть
сказано в возможности формирования поля клиентом.
%
Сервер не должен переносить поле \sstr{onlyReturnExisting}, а также любые другие 
поля, которые он не распознает, в итоговый объект учетной записи. Последнее 
позволяет клиенту определить, поддерживает ли сервер то или иное дополнительное 
поле.

Серверу следует проверить, что адреса в поле \sstr{contact} корректны и 
поддерживаются сервером. Сервер должен проверять адреса схемы \sstr{mailto}. 
Если клиент указывает адрес такой схемы, то он не должен использовать
компонент \sstr{hfields} (см.~\cite{RFC6068}) и не должен указывать более 
одного значения \sstr{addr-spec} в компоненте \sstr{to}. 
%
Если сервер обнаруживает нарушение данных правил, то он должен отклонить адрес
как недействительный. 

Если сервер отклоняет контактный адрес из-за использования неподдерживаемой  
схемы, то он должен возвратить ошибку типа \sstr{unsupportedContact} с 
подробностями проблемы и указанием типов контактных URL, которые сервер считает 
приемлемыми. Если сервер отклоняет контактный адрес из-за использования 
поддерживаемой схемы, но недопустимого значения, он должен возвратить ошибку 
типа \sstr{invalidContact}.  

Если сервер требует от клиента согласия с условиями использования сервиса ACME, 
он должен указать URL-адрес, по которому можно ознакомиться с этими условиями, 
в подполе \sstr{termsOfService} поля \sstr{meta} каталога ресурсов 
(см.~\ref{ACME.Res.Dir}). Сервер должен отклонять запросы на узел 
\code{newAccount}, в которых поле \sstr{termsOfServiceAgreed} либо отсутствует,
либо не принимает значения \code{true}. 
%
Клиенту не следует автоматически соглашаться с условиями обслуживания.
Следует предусмотреть интерактивное взаимодействие с клиентом для согласия с 
условиями.

При успешной обработке запроса на узел \code{newAccount} сервер создает учетную 
запись и сохраняет открытый ключ, который размещен в поле \sstr{jwk} заголовка 
объекта JWS запроса. Открытый ключ будет использоваться в дальнейшем при 
проверке запросов клиента~--- владельца учетной записи. 
%%
Сервер возвращает \pdfexample{ответ}{acme/account2.txt} с кодом статуса 201 
(создан), URL-адресом учетной записи в заголовке \code{Location} и объектом 
учетной записи в теле.
%
URL-адрес учетной записи используется в качестве значения поля \sstr{kid} 
объекта JWS в последующих запросах для данной учетной записи 
(см.~\ref{ACME.Transport.JWS}).
%
URL-адрес учетной записи используется также в запросах на управление этой 
учетной записью, как описано ниже. 

\if0
\begin{codeblock}
HTTP/1.1 201 Created
Content-Type: application/json
Replay-Nonce: D8s4D2mLs8Vn-goWuPQeKA
Link: <https://example.com/acme/directory>;rel="index"
Location: https://example.com/acme/acct/evOfKhNU60wg

{
  "status": "valid",
  "contact": [
    "mailto:cert-admin@example.org",
    "mailto:admin@example.org"
  ],
  "orders": "https://example.com/acme/acct/evOfKhNU60wg/orders"
}
\end{codeblock}
\fi

\subsection{Поиск учетной записи по ключу} \label{ACME.Account.Find}

Если запрос на узел \code{newAccount} признается корректным при проверке на 
открытом ключе, для которого уже зарегистрирована учетная запись, то сервер 
должен возвратить ответ с кодом статуса 200 и указать URL-адрес учетной
записи в заголовке \code{Location}. 
%
В тело ответа сервер включает объект учетной записи в том виде, в котором он 
находился до запроса.
%
Любые поля, которые клиент указывает в объекте своего запроса, должны быть 
проигнорированы сервером. Такая логика позволяет клиенту, у которого есть
ключ учетной записи, но нет соответствующего адреса записи, восстановить адрес.

Если клиент хочет найти URL-адрес существующей учетной записи и не хочет
создавать новую, если она еще не существует, то ему следует отправить 
на узел \code{newAccount} запрос, в объекте которого есть 
поле \sstr{onlyReturnExisting} и оно принимает значение \code{true}.
%
Если клиент отправляет такой запрос, а учетная запись не существует, то сервер
должен возвратить ответ с кодом статуса 400 (неверный запрос) и типом
\sstr{accountDoesNotExist}.

\subsection{Обновление учетной записи}\label{ACME.Account.Upd}

Для обновления учетной записи клиент отправляет 
\pdfexample{запрос}{acme/account3.txt} на URL-адрес учетной записи. В теле 
запроса клиент указывает обновленный объект учетной записи.
%
Сервер должен игнорировать любые обновления полей \sstr{orders}, 
\sstr{termsOfServiceAgreed} (см.~\ref{ACME.Account.ToS} и \sstr{status} 
(за исключением случаев, предусмотренных в~\ref{ACME.Account.Deact}), 
а также любых других полей, которые сервер не распознает. 

Если сервер принимает обновление, то он должен возвратить ответ с кодом статуса 
200 и обновленным объектом учетной записи в теле. 

\if0
Например, чтобы обновить контактную информацию в указанной выше учетной записи, 
клиент может отправить следующий запрос:  

\begin{codeblock}
POST /acme/acct/evOfKhNU60wg HTTP/1.1
Host: example.com
Content-Type: application/jose+json

{
  "protected": base64url({
    "alg": "BIGNS128",
    "kid": "https://example.com/acme/acct/evOfKhNU60wg",
    "nonce": "ax5RnthDqp_Yf4_HZnFLmA",
    "url": "https://example.com/acme/acct/evOfKhNU60wg"
  }),
  "payload": base64url({
    "contact": [
      "mailto:certificates@example.org",
      "mailto:admin@example.org"
    ]
  }),
  "signature": "hDXzvcj8T6fbFbmn...rDzXzzvzpRy64N0o"
}
\end{codeblock}
\fi

\subsection{Изменения условий обслуживания}\label{ACME.Account.ToS} 

Клиент подтверждает согласие с условиями обслуживания сервера через
значение \code{true} в поле \sstr{termsOfServiceAgreed} объекта своей учетной 
записи. 

Если сервер изменил условия обслуживания с момента первоначального согласия 
клиента и не желает обрабатывать определенный запрос без явного согласия с 
новыми условиями, то он должен возвратить \pdfexample{ответ}{acme/account4.txt} 
с кодом статуса 403 (запрещено) и типом ошибки \sstr{userActionRequired}. Этот 
ответ должен включать заголовок \code{Link} с веб-отношением 
\sstr{terms-of-service} и URL-адресом актуальной редакции условий обслуживания.

Подробности проблемы, возвращаемые с ответом, должны также включать поле 
\sstr{instance} с URL-адресом, на который клиент должен направить 
человека-пользователя для получения инструкций о том, как согласиться с 
условиями. 

\if0
\begin{codeblock}
HTTP/1.1 403 Forbidden
Replay-Nonce: T81bdZroZ2ITWSondpTmAw
Link: <https://example.com/acme/directory>;rel="index"
Link: <https://example.com/acme/terms/2017-6-02>;rel="terms-of-service"
Content-Type: application/problem+json
Content-Language: en

{
  "type": "urn:ietf:params:acme:error:userActionRequired",
  "detail": "Terms of service have changed",
  "instance": "https://example.com/acme/agreement/?token=W8Ih3PswD-8"
}
\end{codeblock}
\fi

\subsection{Привязка к внешней учетной записи}\label{ACME.Account.Binding}

Сервер может требовать, чтобы в объекте учетной записи запроса на узел
\code{newAccount} присутствовало поле \sstr{externalAccountBinding} 
или, другими словами, чтобы учетная запись клиента была привязана к внешней 
учетной записи, например, к базе данных клиентов УЦ.

Для привязки к внешней учетной записи УЦ, управляющий сервером ACME, 
должен за пределами ACME предоставить клиенту ключ имитозащиты и его 
идентификатор. Идентификатор ключа должен быть строкой КОИ-7 (ASCII). 
Ключ имитозащиты должен быть представлен кодом base64url. 

Используя ключ имитозащиты, клиент формирует связывающий объект JWS.
Этот объект \pdfexample{указывается}{acme/account5.txt} в поле 
\sstr{externalAccountBinding} родительского объекта JWS. Полезными данными 
связывающего объекта является ключ учетной записи в том виде, в котором он 
задается в поле \sstr{jwk} родительского объекта.
%
Заголовок связывающего объекта должен быть защищен. Его поля должны заполняться
следующим образом:
\begin{itemize}
\item
\str{alg}~--- используемый алгоритм имитозащиты.
Может использоваться алгоритм \algname{hmac-hbelt},
определенный в СТБ 34.101.47. Этому алгоритму назначается 
имя \sstr{HMACBELT};
\item
\sstr{kid}~--- идентификатор ключа имитозащиты;
\item
\sstr{url}~--- повтор одноименного поля родительского JWS.
\end{itemize}
Поле \sstr{nonce} должно отсутствовать.

Поле \sstr{signature} связывающего JWS будет содержать имитовставку открытого
ключа учетной записи, вычисленную на ключе имитозащиты, представленном УЦ.

\if0
\begin{codeblock}
POST /acme/new-account HTTP/1.1
Host: example.com
Content-Type: application/jose+json

{
  "protected": base64url({
    "alg": "BIGNS128",
    "jwk": {...},
    "nonce": "K60BWPrMQG9SDxBDS_xtSw",
    "url": "https://example.com/acme/new-account"
  }),
  "payload": base64url({
    "contact": [
      "mailto:cert-admin@example.org",
      "mailto:admin@example.org"
    ],
    "termsOfServiceAgreed": true,

    "externalAccountBinding": {
      "protected": base64url({
        "alg": "HMACBELT",
        "kid": ...,
        "url": "https://example.com/acme/new-account"
      }),
      "payload": base64url({...}),
      "signature": APqUMnKOoIx0MZ-X7vjTftbHgcKCfvtJ_vVdm_Vlgz0
    }
  }),
  "signature": "5TWiqIYQfIDfALQv...x9C2mg8JGPxl5bI4"
}
\end{codeblock}
\fi

Если сервер требует, чтобы запросы на узел \code{newAccount} содержали поле
\sstr{externalAccountBinding}, он должен указать значение \code{true}
в подполе \sstr{externalAccountRequired} поля \sstr{meta} каталога
ресурсов. Если сервер получает запрос на узел \code{newAccount} без
\sstr{externalAccountBinding}, то ему следует дать ответ с ошибкой типа
\sstr{externalAccountRequired}.

Когда сервер получает запрос на узел \sstr{newAccount} с полем
\sstr{externalAccountBinding}, он решает, проверять ли привязку учетной записи.
Если привязка не проверяется, то сервер не должен включать поле
\sstr{externalAccountBinding} в итоговый объект учетной записи (если таковой
будет сформирован).

Для проверки привязки учетной записи сервер должен выполнить следующие 
действия:

\begin{enumerate}
\item 
Проверить, что в поле указан правильно сформированный объект JWS.
\item
Проверить, что защищаемые поля объекта JWS соответствует вышеуказанным 
правилам.  
\item 
Определить ключ имитозащиты, который соответствует идентификатору 
в поле \sstr{kid}.
\item 
Проверить корректность имитовставки в поле \sstr{signature}. 
\item 
Проверить, что полезные данные объекта JWS повторяют значение поля
\sstr{jwk} родительского JWS.
\end{enumerate}

Если все проверки пройдены успешно и сервер создает новую учетную 
запись, то он может связать эту запись с внешней учетной записью, 
соответствующей ключу имитозащиты.  
%
Возвращаемый сервером объект учетной записи должен включать поле 
\sstr{externalAccountBinding} с тем же значением, что и в запросе. 
%
Если хотя бы одна из проверок не пройдена, сервер 
должен отклонить запрос \code{newAccount}. 

% todo: сервер или УЦ?

\subsection{Обновление ключа учетной записи}\label{ACME.Account.Rollover}

Клиент может изменить ключи, связанные с учетной записью, обратившись с 
запросом на узел \code{keyChange}. Тело \pdfexample{запроса}{acme/account6.txt} 
в этом случае представляет собой подписанные данные в виде объекта JWS.
%
Данные подписываются на действующем личном ключе учетной записи. Подписанные 
данные сами являются объектом JWS (внутренним). Подпись внутреннего объекта JWS 
вырабатывается на новом личном ключе и при этом подписывается объект JSON
со следующими полями:

\begin{itemize}
\item
\sstr{account} (обязательное, строка)~--- URL-адрес учетной записи. 
Адрес должен повторять строку в заголовке \code{Location} ответа с узла 
\code{newAccount} во время создания учетной записи;

\item
\sstr{oldKey} (обязательное, JWK)~--- действующий открытый ключ, повтор
поля \sstr{jwk} из заголовка родительского объекта JWS. 
\end{itemize}

% skip: The signature by the new key covers the account URL and the
%  old key, signifying a request by the new key holder to take over the
%  account from the old key holder.  The signature by the old key covers
%  this request and its signature, and indicates the old key holder's
%  assent to the rollover request.

Внешний объект JWS должен соответствовать стандартным требованиям,
заданным в~\ref{ACME.Transport.JWS}. Внутренний объект JWS также должен 
соответствовать стандартным требованиям со следующими уточнениями:

\begin{itemize}
\item
в поле \sstr{jwk} заголовка должен быть указан новый открытый ключ;
\item
поле \sstr{url} заголовка должно совпадать с одноименным 
полем заголовка внешнего JWS;
\item
в заголовке должно быть опущено поле \sstr{nonce}. 
\end{itemize}

% skip: This transaction has signatures from both the old and new keys so
%  that the server can verify that the holders of the two keys both
%  agree to the change.  The signatures are nested to preserve the
%  property that all signatures on POST messages are signed by exactly
%  one key.  The "inner" JWS effectively represents a request by the
%  holder of the new key to take over the account form the holder of the
%  old key.  The "outer" JWS represents the current account holder's
%  assent to this request.

\if0
\begin{codeblock}
POST /acme/key-change HTTP/1.1
Host: example.com
Content-Type: application/jose+json

{
  "protected": base64url({
    "alg": "BIGNS128",
    "kid": "https://example.com/acme/acct/evOfKhNU60wg",
    "nonce": "S9XaOcxP5McpnTcWPIhYuB",
    "url": "https://example.com/acme/key-change"
  }),
  "payload": base64url({
    "protected": base64url({
      "alg": "BIGNS128",
      "jwk": {...},
      "url": "https://example.com/acme/key-change"
    }),
    "payload": base64url({
      "account": "https://example.com/acme/acct/evOfKhNU60wg",
      "oldKey": {...}
    }),
    "signature": "Xe8B94RD30Azj2ea...8BmZIRtcSKPSd8gU"
  }),
  "signature": "5TWiqIYQfIDfALQv...x9C2mg8JGPxl5bI4"
}
\end{codeblock}
\fi

При получении запроса на узел \code{keyChange} сервер должен выполнить 
следующие шаги в дополнение к обычной проверке объекта JWS:

\begin{enumerate}
\item 
Проверить, соответствует ли запрос учетной записи.
\item 
Проверить, что полезные данные объекта JWS запроса представляют собой правильно 
сформированный объект JWS (внутренний JWS).  
\item 
Проверить, что защищенный заголовок внутреннего JWS имеет поле \sstr{jwk}.
\item
Проверить, что подпись внутреннего JWS признается корректной при проверке на открытом
ключе в поле \sstr{jwk} заголовка.  
\item 
Проверить, что полезные данные внутреннего JWS представляют собой 
правильно сформированный объект JSON с полями \sstr{account} и \sstr{oldKey}.
\item 
Проверить, что поля \sstr{url} в заголовках внутреннего и внешнего JWS 
совпадают.
\item
Проверить, что поле \sstr{account} в полезных данных внутреннего JWS
содержит URL-адрес учетной записи, соответствующий старому ключу 
(т. е. совпадает с полем \sstr{kid} в родительском JWS). 
\item
Проверить, что поле \sstr{oldKey} в полезных данных внутреннего JWS
совпадает с открытым ключом учетной записи.
\item 
Проверить, что не существует учетной записи, открытый ключ которой совпадает с  
ключом в поле \sstr{jwk} заголовка внутреннего JWS.  
\end{enumerate}

Если все проверки пройдены успешно, то сервер обновляет учетную запись, заменяя 
старый открытый ключ учетной записи новым открытым ключом, и возвращает  
ответ с кодом статуса 200. 
%
В противном случае сервер отвечает кодом статуса ошибки и подробностями 
проблемы. 
%
Если обнаружена учетная запись, которая связана с новым открытым ключом, 
серверу следует установить код статуса 409 (конфликт) и указать URL-адрес 
найденной учетной записи в заголовке \code{Location}.

Изменение открытого ключа учетной записи следует организовать так, чтобы
оно не оказывало никакого другого влияния на учетную запись. 
Например, сервер не должен аннулировать отложенные заявки или объекты
авторизации из-за изменения ключа. 

\subsection{Деактивация учетной записи}\label{ACME.Account.Deact}

Клиент может деактивировать учетную запись, отправив подписанное 
\pdfexample{обновление}{acme/account7.txt} учетной записи с полем статуса 
\sstr{deactivated}. Клиент может это сделать, если ключ учетной записи 
скомпрометирован или выведен из эксплуатации.

Сервер должен проверить, что запрос подписан на ключе учетной записи. 
Если сервер принимает запрос на деактивацию, он дает ответ с кодом состояния 
200 и текущим содержимым объекта учетной записи. 

После деактивации учетная запись не может использоваться для выпуска 
сертификатов, ее ресурсы, например, объекты авторизации и заявки, становятся
недоступными. 
%
Если сервер получает запрос на адрес деактивированной учетной записи, 
он должен дать ответ об ошибке с кодом статуса 401 (не авторизовано) и типом 
\sstr{unauthorized}.  

\if0
\begin{codeblock}
POST /acme/acct/evOfKhNU60wg HTTP/1.1
Host: example.com
Content-Type: application/jose+json

{
  "protected": base64url({
    "alg": "BIGNS128",
    "kid": "https://example.com/acme/acct/evOfKhNU60wg",
    "nonce": "ntuJWWSic4WVNSqeUmshgg",
    "url": "https://example.com/acme/acct/evOfKhNU60wg"
  }),
  "payload": base64url({
    "status": "deactivated"
  }),
  "signature": "earzVLd3m5M4xJzR...bVTqn7R08AKOVf3Y"
}
\end{codeblock}
\fi

После деактивации сервер не должен принимать запросы, подписанные на ключе этой 
учетной записи. Сервер должен отменить все ожидающие операции, авторизованные 
ключом этой учетной записи, например, заявки на выпуск сертификатов. 

Сервер может предпринять различные действия в ответ на деактивацию учетной
записи, например, удалить данные, связанные с учетной записью, или
отправить письмо по контактным адресам учетной записи. Сервер не должен отзывать
сертификаты, связанные с деактивированной учетной записью.
%
Повторная активация деактивированной учетной записи в ACME не предусмотрена. 
